%Abstract
\chapter*{Abstract}
\addcontentsline{toc}{chapter}{Abstract}
Crop diseases pose a persistent threat to agricultural productivity in Bangladesh, where rice, potato, and wheat are among the most economically significant crops, yet most existing deep learning-based disease detection systems are trained and evaluated on the PlantVillage dataset, whose clean, laboratory-condition images limit real-world generalizability to natural farm-field conditions. This thesis addresses that gap by constructing an integrated, field-condition dataset of 6,551 images across 15 disease classes for rice, potato, and wheat, assembled from multiple publicly available, locally relevant sources rather than PlantVillage. Using this dataset, three pre-trained CNN backbones -- DenseNet121, MobileNetV2, and EfficientNetB0 -- were evaluated through a progressive pipeline comprising transfer learning with feature extraction, fine-tuning, decision-level fusion, and feature-level fusion. Fine-tuning improved all three backbones over their feature-extraction baselines, with fine-tuned MobileNetV2 achieving the best individual accuracy of 93.65\%, while a decision-level fusion ensemble of the three fine-tuned models underperformed at 91.23\%, likely due to overfitting of its class-specific weighting scheme. In contrast, a feature-level fusion architecture that concatenated the internal deep feature representations of all three fine-tuned backbones achieved the best overall performance, reaching 96.27\% accuracy, 0.97 precision, and 0.96 recall and F1-score on a held-out, unseen test set. These results demonstrate that, for a challenging, locally sourced, multi-crop disease detection problem, combining the internal morphological features of multiple CNN architectures is substantially more effective than combining their final decision outputs, and establish the proposed feature-level fusion model as a robust, field-validated foundation for future real-world crop disease detection systems in Bangladesh.
\clearpage
\setlength{\headheight}{12pt}